\documentclass[border=12pt]{standalone}
\usepackage[utf8]{inputenc}
\usepackage[T1]{fontenc}
\usepackage[spanish,es-noshorthands]{babel}
\usepackage{tikz}
\usetikzlibrary{positioning, shapes.geometric, shapes.symbols, arrows.meta, calc, fit, backgrounds}

% ---- Paleta de colores (misma que los diagramas del TC4) ----
\definecolor{azulPC}{RGB}{33,102,172}      % PC / App Qt
\definecolor{azulPCbg}{RGB}{222,235,247}
\definecolor{verdeESP}{RGB}{35,132,67}     % ESP32
\definecolor{verdeESPbg}{RGB}{225,245,229}
\definecolor{grisRed}{RGB}{110,110,110}    % Red WiFi
\definecolor{grisRedbg}{RGB}{238,238,238}
\definecolor{naranjaHW}{RGB}{217,120,20}   % Hardware
\definecolor{naranjaHWbg}{RGB}{253,237,214}
\definecolor{textoOsc}{RGB}{40,40,40}

\begin{document}
\begin{tikzpicture}[
  font=\sffamily,
  >={Latex[length=2.6mm]},
  paso/.style={
    rectangle, rounded corners=4pt, draw, line width=0.9pt, fill=white,
    align=center, inner sep=5pt, minimum width=44mm, minimum height=10mm,
    text=textoOsc, font=\sffamily\small
  },
  pasoPC/.style={paso, draw=azulPC},
  pasoESP/.style={paso, draw=verdeESP},
  marco/.style={rectangle, rounded corners=8pt, line width=1.1pt, inner sep=7pt},
  hw/.style={
    rectangle, rounded corners=4pt, draw=naranjaHW, fill=naranjaHWbg,
    line width=0.9pt, align=left, inner sep=6pt, text=textoOsc
  },
  nube/.style={
    cloud, cloud puffs=13, cloud puff arc=110, aspect=2.0,
    draw=grisRed, fill=grisRedbg, line width=1pt, align=center,
    inner sep=2pt, minimum width=42mm, minimum height=24mm, text=textoOsc
  },
  etiq/.style={font=\sffamily\footnotesize, align=center, text=textoOsc},
  etiqRet/.style={font=\sffamily\footnotesize, align=center, text=grisRed},
  flecha/.style={->, line width=1.0pt}
]

% ======================= PC: App Qt (columna izquierda) =======================
\node[pasoPC] (desl) {\texttt{deslizadorAngulo} / \texttt{cajaAngulo}\\\scriptsize botones 0° · 90° · 180°};
\node[pasoPC, below=7mm of desl] (sol) {\texttt{solicitarAngulo()}\\\scriptsize máx. 1 mensaje / 50 ms (\texttt{QTimer})};
\node[pasoPC, below=7mm of sol] (env) {\texttt{enviarJson()}\\\scriptsize \texttt{QWebSocket::sendTextMessage}};
\draw[flecha, azulPC] (desl) -- (sol);
\draw[flecha, azulPC] (sol) -- (env);

% ======================= Red WiFi (a la altura de enviarJson) =======================
\node[nube, right=30mm of env] (red) {\textbf{Red WiFi}\\\textbf{SSID GWN571D04}\\[2pt]\footnotesize\texttt{ws://<IP>:81}};

% ======================= ESP32 (columna derecha) =======================
\node[pasoESP, right=40mm of red] (evento) {\texttt{alEventoWebSocket()}\\\scriptsize \texttt{WStype\_TEXT}};
\node[pasoESP, below=7mm of evento] (proc) {\texttt{procesarComando()}\\\scriptsize valida y recorta a 0..180};
\node[pasoESP, below=7mm of proc] (write) {\texttt{servo.write(angulo)}\\\scriptsize PWM 50 Hz, 500–2400\,µs};
\node[pasoESP, below=7mm of write] (ack) {\texttt{avisarAngulo()}\\\scriptsize \texttt{broadcastTXT} a todos};
\draw[flecha, verdeESP] (evento) -- (proc);
\draw[flecha, verdeESP] (proc) -- (write);
\draw[flecha, verdeESP] (write) -- (ack);

% Regreso en la App Qt, a la altura de avisarAngulo()
\node[pasoPC] (interp) at (env |- ack) {\texttt{interpretarEvento()}\\\scriptsize \texttt{etiquetaConfirmado} (ángulo real)};

\begin{scope}[on background layer]
  \node[marco, draw=azulPC, fill=azulPCbg, fit=(desl)(interp),
        label={[font=\sffamily\bfseries, text=azulPC]above:PC — App Qt (cliente)}] (pc) {};
  \node[marco, draw=verdeESP, fill=verdeESPbg, fit=(evento)(ack),
        label={[font=\sffamily\bfseries, text=verdeESP]above:ESP32-S3 — servidor WebSocket :81}] (esp) {};
\end{scope}

% ======================= Ida: comando "servo" =======================
\draw[flecha, azulPC, line width=1.1pt] (env.east) -- node[etiq, above]{\scriptsize WebSocket}
      node[etiq, below]{\scriptsize sobre WiFi} (red.west);
\draw[flecha, verdeESP, line width=1.1pt] (red.east) -- node[etiq, above]{\scriptsize\texttt{\{"tipo":"servo",}}
      node[etiq, below]{\scriptsize\texttt{"angulo":90\}}} (evento.west);

% ======================= Regreso: confirmación =======================
\draw[flecha, grisRed, line width=1.1pt] (ack.west) -- node[etiqRet, above]{\scriptsize\texttt{\{"tipo":"servo\_ack","angulo":90\}}}
      node[etiqRet, below]{\scriptsize mismo WebSocket, sentido contrario} (interp.east);

% ======================= Hardware =======================
\node[hw, right=14mm of write, xshift=7pt] (sg90) {%
  \textbf{Servo SG90}\\[3pt]
  \textbullet~naranja: señal $\leftarrow$ \texttt{GPIO18}\\
  \textbullet~rojo: +5\,V\\
  \textbullet~café: GND (común)};
\draw[flecha, naranjaHW, line width=1.1pt] (write.east) -- node[etiq, above, text=naranjaHW, pos=0.6]{\scriptsize pulso} (sg90.west);

% ======================= Nota al pie =======================
\node[etiq, below=9mm of $(interp.south)!0.5!(ack.south)$, text width=200mm] {%
  \itshape Al conectar, la ESP32 manda \texttt{\{"tipo":"saludo",...,"angulo":n\}}; si el JSON es inválido o falta \texttt{angulo},\\
  responde \texttt{\{"tipo":"error","mensaje":"..."\}} solo a ese cliente.};

\end{tikzpicture}
\end{document}
